\chapter{Parent-origin масс трёх дополнительных рёбер}
\label{ch:v8-baryon-c0-extra-edge-mass-parent-origin}

\section{Три неэквивалентных блока}

После полнографового замыкания вне прямоугольников $B$ и $M$ остаются
ровно три комплексных поля. Их калибровочные типы наследуются от концов:
\begin{equation}
 R_Y=(1,2)_{1/2},\qquad
 R_u=(3,1)_{5/3},\qquad
 R_d=(3,1)_{2/3},
 \label{eq:v8-extra-edge-mass-representations}
\end{equation}
а комплексные размерности равны
\begin{equation}
 (d_Y,d_u,d_d)=(2,3,3).
 \label{eq:v8-extra-edge-mass-dimensions}
\end{equation}
Эти представления попарно неэквивалентны. Поэтому equivariant-блоки между
ними отсутствуют:
\begin{equation}
 \operatorname{Hom}_{G}(R_i,R_j)=0\qquad(i\ne j).
 \label{eq:v8-extra-edge-mass-schur-separation}
\end{equation}
Это различие существенно: диаграмма конечной спектральной тройки фиксирует
представления и допустимые Dirac-связи, но сама по себе не отождествляет
следы разных центральных компонент \cite{v8-extra-edge-mass-krajewski,
v8-extra-edge-mass-cacic}.

\section{Положительная форма из одного оператора}

Для каждого ребра введём самосопряжённый блок
\begin{equation}
 F_e(z_e)=
 \begin{pmatrix}
 0&z_e I_{d_e}\\
 \overline z_e I_{d_e}&0
 \end{pmatrix},
 \qquad e\in\{Y,u,d\}.
 \label{eq:v8-extra-edge-mass-block}
\end{equation}
Тогда один блочно-диагональный оператор
\begin{equation}
 F_z=F_Y\oplus F_u\oplus F_d
 \label{eq:v8-extra-edge-mass-common-operator}
\end{equation}
даёт при физическом полуследе точную положительную форму
\begin{equation}
 \frac12\Tr F_z^2
 =2|z_Y|^2+3|z_u|^2+3|z_d|^2.
 \label{eq:v8-extra-edge-mass-unweighted-half-trace}
\end{equation}
Следовательно, общий операторный носитель действительно порождает все три
необходимых квадратичных слагаемых. Незавешенный матричный след предлагает
отношение
\begin{equation}
 \mu_Y:\mu_u:\mu_d=2:3:3.
 \label{eq:v8-extra-edge-mass-unweighted-ratio}
\end{equation}
Такое появление скалярных связей из спектрального действия согласуется с
общей конструкцией spectral action и её конечными расширениями
\cite{v8-extra-edge-mass-spectral-action,v8-extra-edge-mass-stephan}.

\section{Почему отношение ещё не выведено}

Алгебра трёх блоков сохраняет центральные проекторы $P_Y,P_u,P_d$.
Поэтому вместе с незавешенным следом допустимо всё положительное семейство
\begin{equation}
 G_p=p_YP_Y+p_uP_u+p_dP_d,
 \qquad p_Y,p_u,p_d>0.
 \label{eq:v8-extra-edge-mass-central-metric}
\end{equation}
Оно меняет физическую форму на
\begin{equation}
 \frac12\Tr(G_pF_z^2)
 =2p_Y|z_Y|^2+3p_u|z_u|^2+3p_d|z_d|^2.
 \label{eq:v8-extra-edge-mass-central-family}
\end{equation}
Например, две верные положительные метрики дают
\begin{equation}
 (p_Y,p_u,p_d)=(1,1,1)\mapsto(2,3,3),\qquad
 (1,2,1)\mapsto(2,6,3).
 \label{eq:v8-extra-edge-mass-two-traces}
\end{equation}
Их массовые векторы не пропорциональны. После удаления общего масштаба
остаётся двумерный положительный симплекс относительных весов:
\begin{equation}
 \mathbb P\mathcal C_p\simeq
 \{[p_Y:p_u:p_d]\mid p_e>0\},\qquad
 \dim\mathbb P\mathcal C_p=2.
 \label{eq:v8-extra-edge-mass-projective-simplex}
\end{equation}
Real-удвоение и полуслед умножают все три коэффициента одинаково и этот
симплекс не устраняют. Превратить прямую сумму в один простой фактор мог бы
только новый физический внедиагональный коннектор; но
\eqref{eq:v8-extra-edge-mass-schur-separation} запрещает его при сохранении
текущей gauge-группы.

\section{Калибровочные индексы не являются селектором}

В нормировке $T(\mathbf2)=T(\mathbf3)=1/2$ матрица трёх gauge-индексов,
столбцы которой упорядочены как $(Y,u,d)$, равна
\begin{equation}
 \mathcal I=
 \begin{pmatrix}
 1/2&25/3&4/3\\
 1/2&0&0\\
 0&1/2&1/2
 \end{pmatrix},
 \qquad
 \det\mathcal I=-\frac74,\quad \rank\mathcal I=3.
 \label{eq:v8-extra-edge-mass-gauge-index-matrix}
\end{equation}
Полный ранг означает отсутствие предсказывающего соотношения между тремя
массами: три независимых секторных коэффициента могут воспроизвести любой
массовый вектор. В частности,
\begin{equation}
 \mathcal I^{\mathsf T}
 \begin{pmatrix}0\\4\\6\end{pmatrix}
 =\begin{pmatrix}2\\3\\3\end{pmatrix}.
 \label{eq:v8-extra-edge-mass-gauge-reconstruction}
\end{equation}
Это реконструкция выбранного отношения, а не его вывод; более того, она
требует нулевого $U(1)$-коэффициента. Абсолютный масштаб также несёт
спектральный момент и cutoff-масштаб, которые этот гейт не фиксирует.

Реестры принимают вид
\begin{equation}
 N_{\rm trace\ shape}=5/5,\qquad
 N_{\rm relative\ mass\ origin}=0/5,\qquad
 N_{\rm absolute\ scale\ origin}=0/2.
 \label{eq:v8-extra-edge-mass-ledgers}
\end{equation}

\begin{theorem}[Условный следовой parent трёх масс]
Один самосопряжённый блочный оператор порождает положительные нормы всех
трёх дополнительных рёбер. При незавешенном физическом полуследе их
коэффициенты относятся как $2:3:3$. Однако это отношение является свойством
выбранного представления trace, а не следствием текущей абстрактной
алгебры: её трёхмерный центр оставляет две независимые относительные
координаты.
\end{theorem}

\begin{nogocriterion}[Запрет невыведенного незавешенного следа]
Нельзя объявлять отношение $2:3:3$ предсказанием до выбора единственного
верного trace-state или построения простого factor-parent. Нельзя заменять
эту недостающую операцию подгонкой трёх gauge-секторных коэффициентов:
полноранговая индексная матрица реконструирует, но не выбирает массовый
вектор.
\end{nogocriterion}

\begin{protocol}\begin{enumerate}
\item дополнительных представлений: \textbf{три попарно неэквивалентных};
\item их комплексные размерности: \textbf{$(2,3,3)$};
\item единый блочный оператор: \textbf{построен};
\item положительная квадратичная trace-форма: \textbf{получена};
\item незавешенное отношение: \textbf{$2:3:3$};
\item Real-полуслед меняет отношение: \textbf{нет};
\item размерность центра: \textbf{$3$};
\item относительных trace-свобод: \textbf{две};
\item gauge-index matrix: \textbf{ранг $3$, определитель $-7/4$};
\item gauge-индексы выбирают отношение: \textbf{нет};
\item trace-shape ledger: \textbf{$5/5$};
\item relative-mass-origin ledger: \textbf{$0/5$};
\item absolute-scale-origin ledger: \textbf{$0/2$};
\item полный parent-origin масс: \textbf{не закрыт};
\item следующий гейт: \textbf{селектор центрального trace-симплекса}.
\end{enumerate}\end{protocol}

Машинный сертификат сохранён в
\path{s2t_v8_baryon_c0_singlet_triplet_central_gap_isotypic_channel_extra_edge_mass_parent_origin_gate_results.json}.

\begin{thebibliography}{9}
\bibitem{v8-extra-edge-mass-spectral-action}
A.~H.~Chamseddine, A.~Connes,
\emph{The Spectral Action Principle}, arXiv:hep-th/9606001.
\bibitem{v8-extra-edge-mass-krajewski}
T.~Krajewski,
\emph{Classification of Finite Spectral Triples}, arXiv:hep-th/9701081.
\bibitem{v8-extra-edge-mass-cacic}
B.~\v{C}a\'ci\'c,
\emph{Moduli Spaces of Dirac Operators for Finite Spectral Triples},
arXiv:0902.2068.
\bibitem{v8-extra-edge-mass-stephan}
C.~A.~Stephan,
\emph{New Scalar Fields in Noncommutative Geometry}, arXiv:0901.4676.
\end{thebibliography}